\section{twist-2 算符}
QCD 因子化把结构函数的矩与部分子分布函数的矩联系起来。
以 $F_1$ 为例，其关系为
\be
\int_0^1 dx~x^{n-1} F_1(x,Q^2) = \sum_a C_{1,n}^a(Q^2/\mu^2)
A_n^{a/h}(\mu)
\label{eq:def_mom}
\ee
其中 $A_n^{a/h}(\mu)$ 是某些局域算符（记为 $O_n$）在强子态（记为 $h$）之间
的约化矩阵元，而 $C_{1,n}^a(Q^2/\mu^2)$ 是相应的 Wilson 系数，
可微扰计算。在 DIS 运动学下，式~\eqref{eq:def_mom} 给出的领头贡献
来自具有最低扭度 $\tau = d_O - n = 2$ 的局域算符，
其中 $d_O$ 是局域算符的物理量纲，$n$ 是其自旋。
在部分子模型中，约化矩阵元与 PDF 的矩相联系：
$A_n^{a/h} = \left\langle x^{n-1} \right \rangle_{a/h}$~\cite{Curci:1980uw,Collins:1981uw}。

局域算符 $O_n$ 的形式可以通过要求它们属于 Lorentz 群的不可约表示加以约束。
这保证重正化过程不会产生与同量纲或更低量纲算符的混合。
不过格点 QCD 模拟是在欧几里得框架下进行的，
欧氏格点上的对称群是超立方群 H($4$)。
对称性的降低使算符无法免于复杂的混合模式，
且 $n$ 越大越复杂。完成连续极限后，O($4$) 对称性得以恢复。

构造 O($4$) 对称群的不可约表示是标准内容
（例如参见文献~\cite{hamermesh1962group}）。
对本工作而言，唯一有关的结果是：无迹且对称化的 $n$ 阶张量
构成 O($4$) 的一个不可约表示。

我们限于非极化结构函数，即考虑对强子极化平均后的强子矩阵元。
其他分布的其他矩可用本文所述同一方法研究。
我们进一步把考虑限制在味非单态算符，以避免与胶子算符混合，
但同一方法也适用于单态算符。

我们考虑如下算符集合
\be
O_n^{rs}(x) = O^{rs}_{\mu_1 \cdots \mu_n}(x) =
\psibar^r(x) \gmuopen1 \lrDmu2 \cdots \lrDmuclose{n} \psi^s(x)\,,
\label{eq:t2}
\ee
其中 $\psi^r$ 是不同味的夸克场（$r \neq s$），
协变导数 $D_\mu = \partial_\mu + G_\mu$，$G_\mu$ 为胶子场。
协变导数的对称化组合
$\lrD_\mu = \frac{1}{2}\left({\overset\rightarrow{D}_\mu} -
{\overset\leftarrow{D}_\mu}\right)$
保证在电荷共轭下有确定的变换行为。

在连续情形下，Lorentz 指标对称化的无迹非单态 twist-2
算符~\footnote{无迹算符记作 $\widehat{O}_n$。}
是乘法重正化的。例如在 $\MS$ 方案中，裸算符 $\widehat{O}_{n,{\text{B}}}^{rs}$
按如下方式重正化
\be
\widehat{O}_n^{rs}(x) = Z^{\MS}_n \widehat{O}_{n,{\text{B}}}^{rs}(x)\,,
\ee
其中 $Z^{\MS}_n$ 为重正化常数。强子矩阵元的内禀非微扰性质
表明应当使用格点 QCD。然而，转动群对称性破缺到超立方群 H($4$)
给重正化增添了复杂性。

在超立方格点上，twist-2 算符 Lorentz 指标的选取必须顾及
降低后的超立方对称群 H($4$)。O($4$) 的不可约表示一般会变成 H($4$) 的可约表示，
从而在重正化时诱发不想要的混合~\cite{Kronfeld:1984zv,Martinelli:1987zd,Beccarini:1995iv,Gockeler:1996mu}。
H($4$) 的不可约表示允许与低维算符混合，
并且使同量纲算符间的混合复杂化。
一个著名的复杂化例子是算符 $O_3$。在连续情形，
对 Lorentz 指标对称化并扣除迹之后，重正化是乘法性的（如在 $\MS$ 方案中）。
但在格点上，诸如 $1/a^2 \delta_{\mu_i \mu_j} \cos( a p_{\mu_j})$
这样的项被 H($4$) 对称性允许，即使扣迹也无法去除~\cite{Kronfeld:1984zv}。
这与如下事实有关：$O_3$ 随 Lorentz 指标选取的不同
属于 H($4$) 的不同不可约表示。
指标全相同的 $O_3$（即 $\mu_1 = \mu_2 = \mu_3$）
与低维矢量流混合，产生正比于 $1/a^2$ 的幂次发散；
对于 $\mu_1 \neq \mu_2 = \mu_3$ 的算符，
如不可约表示 $O_{411} - O_{433}$，
重正化是乘法性的，只带有与同量纲另一算符的小混合。\footnote{
    算符 $O_{411} - O_{433}$ 所属的 H($4$) 不可约表示出现不止一次，
    因此属于这些等价表示的算符之间可能发生混合~\cite{Beccarini:1995iv}。
    微扰论表明该混合在数值上很小。
}
注意 $O_{411}$ 本身也受幂次发散影响，
但减除组合 $O_{411} - O_{433}$ 保证将其消除。
另一种选取是 $\mu_1 \neq \mu_2 \neq \mu_3$。此时算符属于 H($4$) 的另一个不可约表示，
乘法重正化且不引入幂次发散。
然而这最后一个例子虽然在重正化意义上最优，
对计算强子矩阵元却远非最优\footnote{为简单起见，式~\eqref{eq:matrix} 中省略了味指标。}
\be
\left \langle h(p)| \widehat{O}_n |h(p) \right\rangle\ =
p_{\mu_1} \cdots p_{\mu_n} A_n^{h}(\mu)\,,
\label{eq:matrix}
\ee
因为它要求至少在 $2$ 个空间方向上有非零外动量。
这会显著恶化数值格点 QCD 模拟中的信噪比。
对 $O_{411} - O_{433}$，至少需要 $1$ 个空间方向的非零外动量，
尽管需要减除幂次发散，仍有望获得较好的信噪比。
指标全相同的算符带有需用复杂非微扰流程减除的幂次发散，并不可行。
遗憾的是，例如 $O_{444}$ 的矩阵元不需要任何外部空间动量，
本应在信噪比上最优。

对更高维算符（$n \ge 4$），混合模式更加棘手。对 $n=4$，
把 $4$ 个指标取全然不同只会得到仅有的 $2$ 个不可约表示：
全对称与全反对称表示。于是需要在全部 $3$ 个空间方向都有非零外动量
才能计算强子矩阵元~\cite{Beccarini:1995iv}。
对 $n > 4$，混合不可避免。上述复杂重正化模式，
加上计算强子矩阵元对外动量的需求，
实际上阻止了 $n>4$ 时部分子分布函数矩的格点计算，
并使 $n=3,4$ 的计算颇具挑战。

文献~\cite{Davoudi:2012ya} 基于连续转动对称性的恢复
提出了绕开重正化问题的有趣想法。
尽管该方法未在数值上继续推进，
它仍是解决 PDF 矩计算中所遇理论挑战的重要第一步。
